
\usepackage{verbatim, multicol, tabularx,hyperref, tikz, enumitem}
\usepackage{amsmath,amsthm, amssymb, stmaryrd, latexsym, bm, listings, qtree}
\usepackage[margin=1in]{geometry}

\lstset{
  extendedchars=\true,
  inputencoding=utf8,
  literate=
  {é}{{\'{e}}}1
  {è}{{\`{e}}}1
  {ê}{{\^{e}}}1
  {ë}{{\¨{e}}}1
  {û}{{\^{u}}}1
  {ù}{{\`{u}}}1
  {â}{{\^{a}}}1
  {à}{{\`{a}}}1
  {î}{{\^{i}}}1
  {ô}{{\^{o}}}1
  {ç}{{\c{c}}}1
  {Ç}{{\c{C}}}1
  {É}{{\'{E}}}1
  {Ê}{{\^{E}}}1
  {À}{{\`{A}}}1
  {Â}{{\^{A}}}1
  {Î}{{\^{I}}}1
  {Ö}{{\"O}}1
  {Ä}{{\"A}}1
  {Ü}{{\"U}}1
  {ö}{{\"o}}1
  {ä}{{\"a}}1
  {ü}{{\"u}}1
  {ß}{{\ss}}1
  ,
  aboveskip=1mm,
  belowskip=1mm,
  showstringspaces=false,
  columns=flexible,
  basicstyle={\scriptsize\ttfamily},
  numbers=left,
  frame=single,
  framextopmargin=0pt,
  framexbottommargin=0pt,
  breaklines=true,
  breakatwhitespace=true,
  keywordstyle=\color{blue},
  identifierstyle=\color{violet},
  stringstyle=\color{teal},
  commentstyle=\color{darkgray}
}

\hypersetup{colorlinks=true,urlcolor=blue}

\headheight = 0.05 in
\headsep = 0.05 in
\parskip = 0.05in
\parindent = 0.0in
\floatsep = 0.05in

\DeclareMathOperator*{\argmin}{arg\!min}
\DeclareMathOperator*{\argmax}{arg\!max}

\title{Temporal-Difference Learning Study Guide (AIMA 22, RLAI 6)}
\author{Artificial Intelligence}
\date{}

\begin{document}
\maketitle

\setcounter{section}{0} % So that next \section is 1
\section{Reinforcement Learning}

\begin{questions}

\question What is reinforcement learning?

\begin{solution}[1in]
Reinforcement learning is learning how to map states to actions by interacting with an environment that provides numerical feedback signals, which we call {\it rewards}.  A reinforcement learning agent tries to maximize long-term reward, not just one-step reward.
\end{solution}

\question What is the {\it reward hypothesis}?

\begin{solution}[1in]
All of what we mean by goals and purposes can be well thought of as the maximization of the expected value of the cumulative sum of a received scalar signal (called reward).
\end{solution}


\question What does ``model-free'' mean in the context of reinforcement learning algorithms that assume there is an underlying MDP model of the environment?

\begin{solution}[1in]
Model-free RL algorithms don't know the underlying MDP, so they learn policies or value functions directly by sampling interactions with the environment.
\end{solution}

\question How many time steps does TD(0) learning use in its target to update to the value of a state?

\begin{solution}[.5in]
1
\end{solution}


\end{questions}

\end{document}
